\subsection{Background (4-6 pages)}
\begin{enumerate}
    \item \textbf{Physics of Radio Recombination Lines} \\
    - What causes them, what do they reveal astronomically
    
    \item \textbf{Previous Surveys}\\
    - Briefly explain WHAM, GDIGS, and SIGGMA Surveys
    
    \item \textbf{Science Impact of a CHORD RRL Survey}\\
    - Explain novelty of a CHORD survey in relation to previous surveys
\end{enumerate}

The background will draw from the \href{https://www.overleaf.com/read/dshhtymzxprq#bbbc27}{RRL Research Primer} document shared at last committee meeting.

\subsection{Theoretical analysis (25 - 40 pages)}
\begin{enumerate}
    \item \textbf{RFI analysis}\\
    - Use the RFI monitor to show which RRL bands are expected to be viable
    - Also consider other transitions
    \item \textbf{CHORD Sensitivity} \\
    - Theoretical CHORD sensitivity based on radiometer equation + Radio background noise.\\
    - All sky sensitivity maps per frequency\\
    - See figure \ref{fig:total_noise}
    \item \textbf{RRL Line Stacking}\\
    - Combining multiple line signals using variance weighted mean\\
    - All sky sensitivity maps based on line stacking strategies\\
    - Analysis that line stacking SNR changes depending on radio background \\
    - See figure \ref{fig:SNR_by_region}
    \item \textbf{Extension CHORD Beam Shapes}\\
    - Analysis of Extension CHORD locations on angular resolution\\
    - See figures \ref{fig:EXT-chord locations} and \ref{fig:EXT-chord location B}
\end{enumerate}
For a nearlty complete version of theoretical analysis see \href{https://www.overleaf.com/project/690bf7f099630fb66b636329}{RRL CHORD Senstivity Doc}

\subsection{Simulation Experiments (5 - 10 pages)}
\begin{enumerate}
    \item \textbf{Galactic plane using SIGGMA}\\
    - The SIGGMA survey provides RRL intensity measurements but sits outside the CHORD observing range, simulate the CHORD response to SIGGMA data as if it were in range. Aka `what would CHORD detect if it could image the same location'.\\
    - Show sensitivity improvements as a function of survey parameters (observing days, lines to include, etc.)\\
    - Example of CHORD applied to SIGGMA see figure \ref{fig:stacked_h2_region}
    \item \textbf{High latitude using WHAM}\\
    - The WHAM survey provides h-alpha map from visible spectrum, which is a reasonable proxy for RRLs\\
    - Same analysis of given WHAM data what would CHORD detect at high latitude vs survey params
\end{enumerate}